% =====================================================================================================
% Additions for current-report.tex. Each block says where it goes. Numbers of the running example come
% from results/xp-main/soybean/taste1/m3-gain/seed3.trace.jsonl (iteration 1 = trace iteration 0) and
% the keep run with the same seed; ICV and partition gains are re-verified by scripts/worked-example.py.
% =====================================================================================================


% -----------------------------------------------------------------------------------------------------
% (A) ABSTRACT: inside \begin{abstract} ... \end{abstract}
% -----------------------------------------------------------------------------------------------------
Interactive pattern mining learns a user's preferences from rankings of a few patterns at a time.
\textsc{DiSPaLe} captures interactions between items by adding, at every iteration, the sub-pattern that best
discriminates the user's ranking as a temporary feature, whose learned weight is then folded into the item
weights. We study whether selecting several complementary discriminant sub-patterns per iteration improves
the learned preferences. We extend \textsc{DiSPaLe} with five selection rules, including a partition-gain
criterion that rewards sub-patterns separating groups the others cannot, and introduce a simulated user whose
preference consists of independent item combinations. We also correct six errors in the published
implementation and verify that undoing them reproduces it exactly. On three datasets, no selection rule
improves the model as long as the sub-pattern features are folded into item weights, since the folded model
cannot represent combinations of items. Keeping the selected sub-patterns as permanent features improves the
learned model significantly, and selecting several of them per iteration adds a further, smaller gain at no
measurable cost.


% -----------------------------------------------------------------------------------------------------
% (B) PRELIMINARIES: after the first sentence of the section ("...shown in Algorithm~1.")
% -----------------------------------------------------------------------------------------------------
\textsc{DiSPaLe} builds on LetSIP~\cite{…}, which follows the same loop with the base features only, and
samples each query with Flexics~\cite{…}, which draws frequent patterns with probability proportional to
the current ranking function.


% -----------------------------------------------------------------------------------------------------
% (C) CORRECTIONS: replace the whole paragraph "Corrections to the interaction loop." with this one
% -----------------------------------------------------------------------------------------------------
\paragraph{Corrections to the interaction loop.}
While re-implementing the extension we found that the publicly available \textsc{DiSPaLe} code does not
execute Algorithm~\ref{alg:dispale} as written. First, the ranking function learned at the end of each
iteration (line~11) was assigned to a local variable that went out of scope, so every query $\mathcal X^t$ was
drawn with the initial weights $w_F^0$ instead of $\phi^{t-1}$: the learned model never influenced which
patterns were shown to the user. We pass $\phi^t$ to the next call of \textsc{SamplePatterns}; the sampler
itself (Flexics) is unchanged and only receives different weights. Second, \textsc{LearnWeights} resumed
SCD from the inner products $w \cdot x$ of the previous iteration, although the dynamic features had been added
and the weights folded since they were computed, so its gradients were wrong from the second iteration on; we
recompute them. Third, \textsc{UpdateWeights} read the learned dynamic weight at the index given by the number
of feature \emph{types} instead of after the base features, and therefore folded the weight of an unrelated
item (item~$1$) instead of that of the sub-pattern. Finally, the code initialises all weights to $1$, not to $w_F^0 = 0$ as
stated in the paper (line~5). With binary item features, the initial score of a pattern is then its length.
We keep this initialisation, since starting from $0$ performs clearly worse (Section~\ref{subsec:fidelity}).

% and add at the end of "Corrections to the sub-pattern search.":
Each of the six corrections can be undone individually. With all of them undone, our implementation
reproduces a run of the published code exactly: the same queries, the same sub-patterns and the same learned
weights at every iteration. Smaller fixes that do not affect our setting (an item index that is wrong when
items are not the first feature type, stored pairs that could only grow by one feature, and two launcher
options) are documented with the code.


% -----------------------------------------------------------------------------------------------------
% (D) MINING: replace "The simplest rule, \emph{Top-$m$}, returns ... add no information." with this
% -----------------------------------------------------------------------------------------------------
Sub-itemsets with the same cover in the query split it in exactly the same way, so they would give the learner
the same feature several times; we therefore keep one candidate per distinct cover, the longest (then the first
found). The first candidate is, by construction, the sub-pattern that the original top-1 rule returns, and our
implementation checks this at every iteration. The simplest rule, \emph{Top-$m$}, returns the $m$
highest-ICV candidates. Even with distinct covers, it ignores how the candidates relate to one another: a
second sub-pattern whose cover lies inside that of the first only splits the same group again.

% and in the paragraph after Algorithm~\ref{alg:partition}, after "greedily adds the candidate with the largest
% gain in PICV", insert:
(ties go to the candidate with the higher ICV)


% -----------------------------------------------------------------------------------------------------
% (E) FEATURE EXPANSION: title, and one sentence at the end of Option 1's first paragraph
% -----------------------------------------------------------------------------------------------------
\subsection{Feature Vector Expansion with Multiple Sub-patterns}
% ...
As in \textsc{DiSPaLe}, the frequency and length descriptors are only added when the base features include
frequency and length; in our experiments the base features are the items, so each sub-pattern contributes its
presence descriptor only.


% -----------------------------------------------------------------------------------------------------
% (F) NEW SUBSECTION: after "Weight Updates under Option 2", before \section{Experimental Evaluation}
% -----------------------------------------------------------------------------------------------------
\subsection{Keeping the Sub-patterns}
\label{subsec:keep}
Folding returns a model over item weights only, which cannot represent a preference for items
\emph{together}, however well the sub-patterns are chosen. We therefore also consider keeping them: every
sub-pattern selected at some iteration becomes a permanent presence feature and is never folded. After
iteration $t$, with $\mathcal K^t$ the set of sub-patterns selected so far, the ranking function is
\begin{equation}\label{eq:keep}
  \phi^t(X) = \phi_{\mathrm{logistic}}\Big(\sum_{i \in X} w_i + \sum_{S \in \mathcal K^t} w_S\,[S \subseteq X],\ A\Big).
\end{equation}
A newly selected sub-pattern starts with weight $0$, all collected pairs are described with the current
features, and \textsc{LearnWeights} learns the item and sub-pattern weights together; \textsc{UpdateWeights}
is skipped. The sampler uses the full model. The feature space grows by at most $m$ per iteration ($47$ kept
sub-patterns after $20$ iterations in the run of Example~\ref{…}).


% -----------------------------------------------------------------------------------------------------
% (G) SETUP: add to the first paragraph
% -----------------------------------------------------------------------------------------------------
As in the published code, SCD performs one sweep over the features per iteration; we also test ten sweeps
(\emph{10sw}). LetSIP is \textsc{DiSPaLe} without dynamic features. A folded model is evaluated on its item
weights, a kept model on its item and sub-pattern weights (Eq.~\eqref{eq:keep}).


% -----------------------------------------------------------------------------------------------------
% (H) CONCLUSIONS: last paragraph
% -----------------------------------------------------------------------------------------------------
Future work should test the keep scheme with real users and with preferences that do not consist of item
pairs, limit the growth of the kept feature set, for instance by pruning sub-patterns whose weight stays near
zero, and revisit the selection rule under this scheme, where covering many candidate combinations appeared to
matter more than complementarity within a single query.


% =====================================================================================================
% RUNNING EXAMPLE: one block per subsection of the methodology, in this order
% =====================================================================================================

% -----------------------------------------------------------------------------------------------------
% (R1) In the combination-user part, after Algorithm~\ref{alg:combo}
% -----------------------------------------------------------------------------------------------------
\paragraph{Running example.}
\label{ex:running}
We follow one iteration of a real run on the soybean dataset with $k = 10$ and a combination user who likes
$\{29, 37\}$ (\texttt{stem\_cankers=absent}, \texttt{external\_decay=absent}) and $\{33, 45\}$
(\texttt{canker\_lesion=dna}, \texttt{sclerotia=absent}). Table~\ref{tab:ex_query} shows the first query in
the order the user ranks it. The three patterns that contain $\{33, 45\}$ score $1$ plus a small
surprisingness term and come first; the seven others contain no liked pair and are ordered by surprisingness
alone. Since all weights start at $1$, the model score $w \cdot x$ of every pattern is its length.

\begin{table}[t]
\centering
\caption{Running example: the first query, ranked by the user (rank $0$ preferred). \#: position in the
sampled order; liked pair in bold.}
\label{tab:ex_query}
\small
\setlength{\tabcolsep}{4pt}
\begin{tabular}{rrlrr}
\toprule
Rank & \# & Pattern & $w \cdot x$ & $u(X)$ \\
\midrule
0 & 3 & $\{8,11,27,29,\mathbf{33},40,42,\mathbf{45}\}$ & 8 & 1.0008 \\
1 & 6 & $\{19,\mathbf{33},42,\mathbf{45},47\}$          & 5 & 1.0003 \\
2 & 7 & $\{7,13,18,27,\mathbf{33},\mathbf{45}\}$        & 6 & 1.0002 \\
3 & 5 & $\{11,26,28,32,40\}$                            & 5 & 0.0012 \\
4 & 4 & $\{11,13,28,35,40,42,45\}$                      & 7 & 0.0007 \\
5 & 1 & $\{4,7,23,26,40,45\}$                           & 6 & 0.0004 \\
6 & 9 & $\{10,20,35,40,42\}$                            & 5 & 0.0003 \\
7 & 8 & $\{11,13,23,26,34,40,47\}$                      & 7 & 0.0003 \\
8 & 2 & $\{14,18,27,37,40,42,47\}$                      & 7 & 0.0003 \\
9 & 0 & $\{7,11,23,32,37,40,45\}$                       & 7 & 0.0000 \\
\bottomrule
\end{tabular}
\end{table}

% -----------------------------------------------------------------------------------------------------
% (R2) In "Mining multiple discriminating sub-patterns", right after Eq.~\eqref{eq:icv} and its "where ..."
% -----------------------------------------------------------------------------------------------------
\paragraph{Running example (cont.).}
In Table~\ref{tab:ex_query}, $\mu = 4.5$. The liked pair $Y = \{33, 45\}$ is contained in the patterns of
ranks $0, 1, 2$, so
\begin{equation*}
  \ICV(Y) = 3\,(4.5 - 1)^2 + 7\,(4.5 - 6)^2 = 36.75 + 15.75 = 52.5 ,
\end{equation*}
the highest of the $39$ distinct covers with a positive ICV. Next come $\{33, 42, 45\}$ (ranks $0, 1$) and
$\{37, 40\}$ (ranks $8, 9$), both with ICV $40.0$.

% -----------------------------------------------------------------------------------------------------
% (R3) After the paragraph that follows Algorithm~\ref{alg:partition}
% -----------------------------------------------------------------------------------------------------
\paragraph{Running example (cont.).}
With $m = 3$, \textsc{PartitionGain} starts from $\{33, 45\}$ ($\PICV = 52.5$; $\TSS = 82.5$). Although
$\{33, 42, 45\}$ has the second-highest ICV, its cover lies inside that of $\{33, 45\}$ and only splits the
preferred group again. $\{37, 40\}$ instead separates the two least preferred patterns, raising PICV to $70.0$
(it ties with $\{11, 28, 40\}$, which has the lower ICV). The third step adds $\{11, 28, 40\}$ (ranks $3, 4$;
$\PICV = 77.5$). The three sub-patterns divide the query into four groups, the preferred ranks $0$--$2$, the
least preferred $8$--$9$, ranks $3$--$4$ and ranks $5$--$7$, which together explain $77.5$ of $82.5$.

% -----------------------------------------------------------------------------------------------------
% (R4) In Option 1, after Eq.~\eqref{eq:patterns-fine}
% -----------------------------------------------------------------------------------------------------
\paragraph{Running example (cont.).}
With the three selected sub-patterns and item features only, each pattern of Table~\ref{tab:ex_query} gets
three presence descriptors: $(1, 0, 0)$ for the preferred pattern of rank $0$, which contains $\{33, 45\}$, and
$(0, 1, 0)$ for the pattern of rank $9$, which contains $\{37, 40\}$. In the training pair formed by these two
patterns, the dynamic part of the difference vector is therefore $(1, -1, 0)$.

% -----------------------------------------------------------------------------------------------------
% (R5) In "Weight Updates under Option 1", after Eq.~\eqref{eq:fine-lin-item}/(eq:fine-exp-item)
% -----------------------------------------------------------------------------------------------------
\paragraph{Running example (cont.).}
After learning on this first query ($\eta = 0.17$, linear update), the weight of item $33$ has grown from $1$ to
$2.16$, but the weight of the liked pair $\{33, 45\}$ is still exactly $0$. In this query, item $33$ occurs in
exactly the same three patterns as $\{33, 45\}$, and SCD updates the item weights before the dynamic ones, so
item $33$ absorbs the ranking signal first. Only $\{37, 40\}$ receives a non-zero
weight, $-0.58$. Folding then multiplies the weights of items $37$ and $40$ by
$1 + 0.17 \cdot (-0.58 / 3) = 0.967$ (item $40$: $1.0 \to 0.967$), and leaves items $33$ and $45$ unchanged.
The preference for $33$ and $45$ \emph{together} is not transferred; it is only partially captured by the
weight of item $33$ alone, and the pair must be rediscovered in a later query.

% -----------------------------------------------------------------------------------------------------
% (R6) At the end of the new subsection (F) "Keeping the Sub-patterns"
% -----------------------------------------------------------------------------------------------------
\paragraph{Running example (cont.).}
In the run with the same seed and the keep scheme, both liked pairs stay in the model. The weight of
$\{29, 37\}$ grows from $0.26$ after iteration $2$ to $1.16$ after iteration $20$, the largest of the $47$ kept
sub-patterns, and that of $\{33, 45\}$ from $0$ to $0.72$.

% -----------------------------------------------------------------------------------------------------
% (R7) Optional: iterations 2 and 3 in the appendix, as table* (two-column width), from
%      docs/report-examples.tex (Table "tab:ex_iter23"); change \begin{table} to \begin{table*}.
% -----------------------------------------------------------------------------------------------------
